\section{Introduction}
\label{sec:introduction}

Retrieval-augmented generation (RAG) systems are increasingly deployed in time-sensitive domains such as financial news analysis, medical literature triage, and breaking-event summarization, where the factual landscape shifts on the scale of days or weeks. Yet the benchmarks used to evaluate these systems---TriviaQA~\citep{joshi2017triviaqa}, Natural Questions~\citep{kwiatkowski2019nq}, HotpotQA~\citep{yang2018hotpotqa}, and BEIR~\citep{thakur2021beir} among them---were constructed from static Wikipedia snapshots and web crawls captured years before the models being evaluated. When a benchmark's knowledge cutoff predates the model's training cutoff, leakage is all but guaranteed: a model that has memorized benchmark answers during pretraining will appear to perform retrieval-augmented reasoning while doing nothing of the sort. Beyond contamination, static benchmarks cannot measure how quickly citation faithfulness decays as indexed corpora age, nor can they detect rank reversals that emerge when a strong model's parametric memory becomes stale relative to a weaker but fresher retriever. As an illustrative example, ColBERT+GPT-4o achieves 80.1 F1 on pre-cutoff (Static) content yet degrades to 71.3 F1 on post-cutoff (Fresh) content---an 8.8-point drop---while DPR+GPT-4o degrades by 13.4 F1 points over the same transition, a rank-altering gap invisible to any frozen benchmark.

Prior work has addressed evaluation quality along several orthogonal axes. Dynamic benchmarking efforts such as Dynabench~\citep{dynabench2021} advocate adversarial human-in-the-loop data collection, but do not track temporal drift in retrieval corpora. FreshQA~\citep{vu2023freshqa} targets world-knowledge freshness for closed-book QA, but does not study retrieval-grounded systems or measure citation faithfulness over time. BEIR~\citep{thakur2021beir} provides a heterogeneous retrieval benchmark, yet its splits are fixed and researchers routinely evaluate on data that postdates their systems' training. Citation evaluation work~\citep{gao2023enabling,min2023factscore} focuses on whether generated claims are grounded in retrieved passages, but does not track how citation support rates evolve as corpora become stale. \method{} addresses all three gaps simultaneously: a rolling monthly benchmark that (i) sources documents and questions exclusively from content published after the models' training cutoffs, (ii) tracks citation precision and recall alongside answer F1 as the benchmark ages, and (iii) detects temporal contamination via n-gram overlap between benchmark questions and model training corpora.

This paper makes the following contributions:
\begin{itemize}
  \item \textbf{Rolling benchmark protocol.} We introduce a principled pipeline for constructing monthly evaluation slices from Reuters, BBC, and AP news feeds and Wikipedia recent-edit streams, with automated temporal contamination detection that filters any question whose answer span appears verbatim in publicly available pretraining data.
  \item \textbf{Temporal degradation analysis.} Characterizing how answer F1 and citation precision degrade as a function of months elapsed since the training cutoff, we find that Fresh F1 falls to the 53--61 range at 12 months post-cutoff (from Static baselines of 74--80), a drop of 8--14 F1 points. The freshness composite score $\mathcal{F}$ correlates with human-judged answer currency at $\rho{=}0.81$, outperforming raw F1 ($\rho{=}0.63$) and citation precision alone ($\rho{=}0.71$).
  \item \textbf{Rank-instability and contamination findings.} System rankings on fresh monthly slices disagree with static benchmark rankings 38\% of the time. Disabling contamination detection inflates apparent Fresh F1 by $+4.9$ points (95\% CI: 3.9--5.5), of which 61\% reflects memorized parametric answers rather than genuine retrieval---confirming that static leaderboards systematically overstate real-world performance.
\end{itemize}
